Leaderboards rank models long after their benchmarks have stopped measuring real differences. A common belief is that the shape of a benchmark's item-difficulty distribution decides when this happens: uniform benchmarks should saturate near 85\% accuracy, heavy-tailed ones should last until 95\%, a benchmark's useful lifetime should track the variance of its item difficulties, and at saturation evaluation-seed noise should outweigh capability in rankings. We test these claims with a matched-variance item response theory (IRT) analysis, per-item responses of 220 \leaderboard models on 62 tasks, real items resampled into prescribed difficulty shapes, and 225{,}000 option-order-permuted evaluations of 15 open LLMs. Shape barely matters. With variance held fixed, all seven shapes lose discrimination within about 1.5 points of each other (95.2--96.6\% accuracy for 500 items), and the hard-tailed shape fails first, not last. Real items reproduce this null. The saturation accuracy is instead set by the number of items (80--84\% at 100 items, 96\% at 500, 99\% at 2000) and by how tightly the competing models cluster. Difficulty variance predicts lifetime only among equal-size benchmarks ($r{=}0.92$ in simulation at $N{\ge}500$), and not across real tasks ($r{=}{-}0.06$, 95\% CI $[-0.38, 0.33]$). Seed noise is about 0.40$\times$ the binomial item-sampling noise. It explains at most 18\% of score variance on 500-item benchmarks and dominates only for near-tied models on about 150 items at $\ge$99\% accuracy. Benchmark builders should spend items near the frontier and keep test sets large; difficulty-histogram shape is a second-order concern.
